% Lecture notes: Week 11, Lecture 3 -- Solving Statically Indeterminate Problems
% To hide this lecture from the website, add a line containing only:  % draft
% To set the date shown on the website, add a line like:  % posted: 2026-09-02
\documentclass[11pt]{article}
\usepackage{mech2030notes}
\begin{document}

% ##################################################################
%  WEEK 11, LECTURE 3
% ##################################################################
\lecture{11}{3}{Solving Statically Indeterminate Problems}

\section*{Example}

\noindent
\begin{minipage}[c]{0.55\textwidth}
\centering
\begin{tikzpicture}
  % scale: 1 cm per m
  \lwall{0}{-0.9}{0.9}
  \rwall{5}{-0.9}{0.9}
  \draw[beam] (0,-0.3) rectangle (5,0.3);
  \node[left] at (-0.4,0) {$A$};
  \node[right] at (5.4,0) {$B$};
  \node[pin] at (2,0) {}; \node[below] at (2,-0.02) {$C$};
  \draw[force] (2.05,0) -- (3.1,0);
  \node[above] at (2.6,0.3) {$\SI{500}{N}$};
  \draw[dim] (0,-1.3) -- (2,-1.3) node[midway, fill=white] {$\SI{2}{m}$};
  \draw[dim] (2,-1.3) -- (5,-1.3) node[midway, fill=white] {$\SI{3}{m}$};
\end{tikzpicture}
\end{minipage}%
\hfill
\begin{minipage}[c]{0.42\textwidth}
A bar is fixed to walls at both ends and loaded axially by $\SI{500}{N}$ at $C$.

\smallskip
\textbf{Q:} Find the reactions at the walls.
\end{minipage}

\subsection*{FBD}

\begin{center}
\begin{tikzpicture}
  \draw[thick] (0,0) -- (5,0);
  \node[pin] at (0,0) {}; \node[pin] at (2,0) {}; \node[pin] at (5,0) {};
  \node[below] at (2,-0.05) {$C$};
  \draw[force] (-1.2,0) node[left] {$A_x$} -- (-0.08,0);
  \draw[force] (2,0.3) -- (3.0,0.3) node[right] {$\SI{500}{N}$};
  \draw[force] (5.08,0) -- (6.2,0) node[right] {$B_x$};
\end{tikzpicture}
\end{center}

\[
  \textstyle\sum F_x = 0 \quad\rightarrow\quad \boxed{A_x + \SI{500}{N} + B_x = 0}
\]
\begin{itemize}[label=\then]
  \item 2 unknown reactions.
  \item 1 statics equation ($\sum F_x = 0$).
  \begin{itemize}[label=\then]
    \item \emph{Statically indeterminate!}
  \end{itemize}
\end{itemize}

\noindent
$\rightarrow$ Introduce an \emph{additional} governing equation:
\[
  \boxed{\delta_B = 0} \qquad \text{(no displacement at the wall)}
\]
\begin{itemize}[label=\then]
  \item Now we have 2 unknowns with 2 equations (solvable!).
  \item An extra equation involving displacement is called a
        ``\emph{compatibility condition}.''
\end{itemize}

\subsection*{Superposition}

\noindent
$\rightarrow$ Use \emph{superposition} to find an expression for $\delta_B$. Remove the wall
at $B$ and replace it with the (unknown) reaction force $B_x$:

\begin{center}
\begin{tikzpicture}[baseline=(b)]
  \coordinate (b) at (0,0);
  \lwall{0}{-0.9}{0.9}
  \rwall{5}{-0.9}{0.9}
  \draw[beam] (0,-0.3) rectangle (5,0.3);
  \node[pin] at (2,0) {};
  \draw[force] (2.05,0) -- (3.1,0);
  \node[above] at (2.6,0.3) {$\SI{500}{N}$};
\end{tikzpicture}

\vspace{1em}
{\Large $=$}\quad
\begin{tikzpicture}[baseline=(b), scale=0.85]
  \coordinate (b) at (0,0);
  \lwall{0}{-0.9}{0.9}
  \draw[beam] (0,-0.3) rectangle (5,0.3);
  \node[pin] at (2,0) {};
  \draw[force] (2.05,0) -- (3.1,0);
  \node[above] at (2.6,0.3) {$\SI{500}{N}$};
  \node[pin] at (5,0) {}; \node[above right] at (5,0.3) {$B$};
  \node at (2.5,-1.0) {\circnum{I}};
  \draw[dim] (0,-1.6) -- (2,-1.6) node[midway, fill=white] {$\SI{2}{m}$};
\end{tikzpicture}
\quad{\Large $+$}\quad
\begin{tikzpicture}[baseline=(b), scale=0.85]
  \coordinate (b) at (0,0);
  \lwall{0}{-0.9}{0.9}
  \draw[beam] (0,-0.3) rectangle (5,0.3);
  \node[pin] at (5,0) {}; \node[above right] at (5,0.3) {$B$};
  \draw[force] (5.05,0) -- (6.1,0) node[below] {$B_x$};
  \node at (2.5,-1.0) {\circnum{II}};
  \draw[dim] (0,-1.6) -- (5,-1.6) node[midway, fill=white] {$\SI{5}{m}$};
\end{tikzpicture}
\end{center}

\noindent
Using $\delta = \frac{NL}{EA}$ for each case (only the loaded length stretches in case
\circnum{I}):
\[
  \delta_B = \delta_B^{\text{I}} + \delta_B^{\text{II}}
  = \frac{(\SI{500}{N})(\SI{2}{m})}{EA} + \frac{B_x\,(\SI{5}{m})}{EA}
\]

\noindent
$\hookrightarrow$ Also, from the compatibility condition:
\[
  \delta_B = 0 = \frac{(\SI{500}{N})(\SI{2}{m})}{EA} + \frac{B_x\,(\SI{5}{m})}{EA}
  \qquad (EA \text{ cancels})
\]
\[
  B_x = -(\SI{500}{N})\left(\frac{\SI{2}{m}}{\SI{5}{m}}\right)
  \qquad\Rightarrow\qquad
  \fbox{\fbox{$B_x = -\SI{200}{N}$}}
\]

\noindent
$\rightarrow$ From the global FBD:
\[
  A_x = -B_x - \SI{500}{N} = \SI{200}{N} - \SI{500}{N}
  \qquad\Rightarrow\qquad
  \fbox{\fbox{$A_x = -\SI{300}{N}$}}
\]

\subsection*{Summary}

\noindent
We can solve statically indeterminate problems by:
\begin{enumerate}
  \item Introducing ``\emph{compatibility conditions}'' involving \emph{deformation} of the
        body as \emph{additional equations}.
  \item Using \emph{superposition} to find displacement expressions, and solving for the
        unknown reactions.
\end{enumerate}

\end{document}
